%! Characteristics of Rational Functions — Unit 4, Lesson 4.0 — Homework (Answer Key)
\documentclass[10pt]{article}
\usepackage{algebra2-article}
\usepackage{algebra2-key}   % pulls in -boxes; do NOT also load -boxes
\newcommand{\TallMath}[1]{$\displaystyle #1\rule[-1.4em]{0pt}{3.2em}$}
\newcommand{\lgrid}[4]{%
  \draw[step=1, linegray!40, very thin] (#1,#3) grid (#2,#4);
  \draw[->, semithick, charcoal] (#1,0)--(#2,0) node[below right, font=\scriptsize]{$x$};
  \draw[->, semithick, charcoal] (0,#3)--(0,#4) node[left, font=\scriptsize]{$y$};}

\begin{document}

\pageheader{Unit 4, Lesson 4.0}{Homework --- Answer Key}
\namedateperiod

\vspace{0.10in}

\begin{notesbox}{Practice}
Show how you read each feature. Write every asymptote as an \textbf{equation}.

\begin{enumerate}[leftmargin=1.9em, itemsep=12pt]
  \item \textbf{From the equation only.} Factor where you can, then fill in each row. Write ``none''
        when there is nothing to report.
    \begin{center}
    \renewcommand{\arraystretch}{1.6}
    \begin{tabularx}{\linewidth}{@{}l X X X X@{}}
    \hline
    \textbf{Function} & \textbf{Excluded value(s)} & \textbf{Vert.\ asymptote(s)} & \textbf{Hole} & \textbf{Horiz.\ asymptote} \\
    \hline
    \TallMath{\frac{5}{x+3}} & \ans{$x=-3$} & \ans{$x=-3$} & \ans{none} & \ans{$y=0$} \\
    \TallMath{\frac{4x}{x-6}} & \ans{$x=6$} & \ans{$x=6$} & \ans{none} & \ans{$y=4$} \\
    \TallMath{\frac{x^2-16}{x+4}} & \ans{$x=-4$} & \ans{none} & \ans{$(-4,-8)$} & \ans{none} \\
    \TallMath{\frac{x-2}{x^2-x-2}} & \ans{$x=2,\,-1$} & \ans{$x=-1$} & \ans{$\left(2,\tfrac13\right)$} & \ans{$y=0$} \\
    \hline
    \end{tabularx}
    \end{center}
    {\small \textbf{Work:} $\dfrac{x^2-16}{x+4}=\dfrac{(x-4)(x+4)}{x+4}=x-4$ ($x\neq-4$), so the hole
    is at $(-4,-8)$ and the top-heavy degrees ($2>1$) leave no horizontal asymptote.
    $\dfrac{x-2}{x^2-x-2}=\dfrac{x-2}{(x-2)(x+1)}=\dfrac{1}{x+1}$ ($x\neq2$): $(x-2)$ cancels
    $\Rightarrow$ hole at $\left(2,\tfrac13\right)$; $(x+1)$ stays $\Rightarrow$ asymptote $x=-1$.}

  \item \textbf{A graph with two asymptotes.} The graph is $m(x) = \dfrac{x-4}{x-2}$.
    \begin{center}
    \begin{tikzpicture}[scale=0.44]
      \lgrid{-4}{8}{-4}{6}
      \draw[navy, dashed, semithick] (2,-4)--(2,6);
      \draw[navy, dashed, semithick] (-4,1)--(8,1);
      \begin{scope}
        \clip (-4,-4) rectangle (8,6);
        \draw[very thick, forest, domain=-4:1.6, samples=90, smooth]
          plot (\x,{((\x)-4)/((\x)-2)});
        \draw[very thick, forest, domain=2.4:8, samples=90, smooth]
          plot (\x,{((\x)-4)/((\x)-2)});
      \end{scope}
      \fill[charcoal] (4,0) circle (3pt) node[below right, font=\scriptsize]{$(4,0)$};
      \fill[charcoal] (0,2) circle (3pt) node[left, font=\scriptsize]{$(0,2)$};
    \end{tikzpicture}
    \end{center}
    Vertical asymptote: \ans{$x=2$} \quad Horizontal asymptote: \ans{$y=1$} \\[3pt]
    Domain: \ans{$(-\infty,2)\cup(2,\infty)$} \quad Range: \ans{$(-\infty,1)\cup(1,\infty)$} \\[3pt]
    $x$-intercept: \ans{$(4,0)$} \quad $y$-intercept: \ans{$(0,2)$} \\[3pt]
    Increasing on: \ans{$(-\infty,2)$} \; and \; \ans{$(2,\infty)$} \\[3pt]
    End behavior: as $x\to\pm\infty$, $m(x)\to$ \ans{1}

  \item \textbf{A graph with a hole.} The graph is $n(x) = \dfrac{x^2-x-6}{x-3}$.
    \begin{center}
    \begin{tikzpicture}[scale=0.46]
      \lgrid{-5}{5}{-3}{7}
      \begin{scope}
        \clip (-5,-3) rectangle (5,7);
        \draw[very thick, forest] (-5,-3)--(5,7);
      \end{scope}
      \draw[forest, very thick, fill=white] (3,5) circle (4.2pt);
    \end{tikzpicture}
    \end{center}
    Factor and simplify: $n(x)=$ \ans{$\frac{(x-3)(x+2)}{x-3}=x+2$}, \; for $x\neq$ \ans{3} \\[3pt]
    The hole is at \ans{$(3,5)$} \quad Vertical asymptote: \ans{none} \quad Horizontal asymptote:
    \ans{none} \\[3pt]
    Domain: \ans{$(-\infty,3)\cup(3,\infty)$} \quad $x$-intercept: \ans{$(-2,0)$} \quad
    $y$-intercept: \ans{$(0,2)$}

  \item \textbf{Explain.} A graph may \emph{cross} its horizontal asymptote, but it can \emph{never}
        cross a vertical asymptote. Explain the difference.
        \ansline{A vertical asymptote sits at an \emph{excluded input}. There is no point on the graph at}
        \ansline{that $x$-value at all, so there is nothing there to cross --- the graph can only run alongside it.}
        \ansline{A horizontal asymptote restricts nothing. It only describes what the \emph{outputs} do at the far ends, so in the middle of the graph an output is free to equal that value, and the curve can pass through the line.}
\end{enumerate}
\end{notesbox}

\begin{extensionbox}
\textbf{Reading a model in context.} After a patient takes one dose of a medication, the concentration
in the bloodstream (in mg/L) $t$ hours later is modeled by
\[
  C(t) = \frac{5t}{t^2+1}, \qquad t \ge 0 .
\]
\begin{center}
\renewcommand{\arraystretch}{1.25}
\begin{tabular}{|c|c|c|c|c|c|c|}
\hline
$t$ (hours) & $0$ & $0.5$ & $1$ & $2$ & $4$ & $8$ \\ \hline
$C(t)$ (mg/L) & $0$ & $2$ & $2.5$ & $2$ & $\approx 1.18$ & $\approx 0.62$ \\ \hline
\end{tabular}
\end{center}
\begin{enumerate}[label=(\alph*), leftmargin=1.8em, itemsep=5pt]
  \item Compare the degrees of the numerator and denominator to find the horizontal asymptote.
        \ans{Numerator degree $1$, denominator degree $2$; bottom-heavy ($n<m$), so the HA is $y=0$.}
  \item What does that horizontal asymptote \emph{mean} about the medication over a long time?
        \ans{As hours pass the concentration drops toward $0$ --- the body clears the drug out. It gets arbitrarily small without ever becoming exactly $0$ in the model.}
  \item From the table, when is the concentration highest, and what is that concentration?
        \ans{At $t=1$ hour, the concentration peaks at $2.5$ mg/L (an absolute maximum on $t\ge0$).}
  \item Notice that $C(0)=0$ --- the graph \emph{sits on} its horizontal asymptote at the start. Does
        that contradict anything from today's lesson?
        \ans{No. A horizontal asymptote governs only the \emph{far ends}; a graph may touch or cross it elsewhere. Only a \emph{vertical} asymptote can never be touched.}
  \item The denominator $t^2+1$ is \emph{never} zero, so this rational function has \textbf{no}
        vertical asymptote and no excluded values at all. Why, then, is the domain written $t\ge 0$?
        \ansline{The restriction comes from the \emph{context}, not the algebra: $t$ is hours since the dose, and}
        \ansline{negative time is meaningless here. A rational function is not required to have a vertical asymptote --- only a denominator that can equal zero produces one.}
\end{enumerate}
\end{extensionbox}

\begin{spiralbox}
\textbf{Coming next --- Lesson 4.1: Simplifying Rational Expressions \& Domain Restrictions.} Today
you \emph{read} holes and asymptotes off a graph. Next you will do the algebra that produces them:
factor a rational expression, cancel common factors, and list the restrictions --- always from the
\textbf{original} denominator. That last habit is the one that will save you in Lesson~4.6, when a
``solution'' turns out to be an excluded value in disguise.
\end{spiralbox}

\begin{teachernote}
\textbf{Item 1} is the lesson's transfer task: rows 3 and 4 cannot be answered without factoring
first, and both are designed to catch the habit of announcing a vertical asymptote at every excluded
value. Row 4 is the hardest --- two excluded values, only one asymptote --- and the hole's $y$-value
$\tfrac13$ must come from the \emph{simplified} expression $\frac{1}{x+1}$, not the original.
\textbf{Item 3} is the mirror of Activity Tier~A: the graph is a plain line, so students who wrote
``vertical asymptote'' anywhere in it need the cancel test re-taught before 4.1. \textbf{Item 4} is
the conceptual item worth the most discussion; accept any answer that distinguishes a restriction on
\emph{inputs} from a description of \emph{outputs}.

\textbf{Extension} part (e) is deliberately unsettling and worth two minutes of class time: a rational
function with \emph{no} vertical asymptote at all. It blocks the overgeneralization ``every rational
function has a vertical asymptote,'' which otherwise resurfaces in 4.5.
\end{teachernote}

\end{document}
